\documentclass[11pt,letterpaper,oneside,reqno]{amsart}
\usepackage[T1]{fontenc}
\usepackage{lmodern}
\usepackage{amssymb}
\usepackage[margin=1.12in]{geometry}
\usepackage{microtype,booktabs}
\setlength{\emergencystretch}{1.5em}
\usepackage[hidelinks]{hyperref}
\hypersetup{pdfauthor={},pdftitle={An explicit power saving for planar unit distances}}
\pagestyle{plain}
\raggedbottom

\newtheorem{mainthm}{Theorem}
\renewcommand{\themainthm}{\Alph{mainthm}}
\newtheorem{lemma}{Lemma}[section]
\theoremstyle{definition}
\newtheorem{definition}[lemma]{Definition}
\title{An explicit power saving for planar unit distances}
\date{}
\newcommand{\eventline}{Round 1 Mathathon \enspace$\cdot$\enspace Geometry
  \enspace$\cdot$\enspace Problem 46}

\newif\ifshowteam
\showteamtrue
\ifshowteam
  \author[]{Deep Blue}
  \dedicatory{\normalfont
    Shubhrajit Bhattacharya and Sofia de la Cerda, University of Chicago\\
    Ritesh Goenka and Dmitry Tsarev, University of Oxford\\[.55em]
    \eventline\\[.55em]
    \url{https://github.com/shubhrajitbhattacharya/unit_distance_formalization}\\
    \href{https://shubhrajit.com/p-e7eddfbaa057d78d86c1d003/}{Interactive proof blueprint}}
\else
  \dedicatory{\normalfont\eventline}
\fi

\begin{document}
\begin{abstract}
OpenAI proved that the number of unit distances among \(n\) points of
the Euclidean plane is \(O(n^{4/3-\delta})\) for some \(\delta>0\).
We explain the prediction argument through Proposition~4.4 of their
paper and track the size of a finite test for the circle equation.
This gives an explicit saving, for which the simpler value
\(2^{-2^{2115}}\) also suffices. In the large-variation case, the
height of one stationary coordinate on an \(8\times4\) grid is
enough to complete the count. We give the key estimates and record
the scope and provenance of the Lean verification.
\end{abstract}
\maketitle
\thispagestyle{plain}
\section{The unit distance problem}
Let \(u(n)\) be the largest number of unordered pairs at Euclidean distance
one among \(n\) distinct points in the plane. Erd\H{o}s's lattice
construction gives \(u(n)\ge n^{1+c/\log\log n}\) for an absolute \(c>0\)
and arbitrarily large \(n\) \cite{E}. Spencer, Szemer\'edi and Trotter
proved \(u(n)=O(n^{4/3})\) \cite{SST}. In May 2026 OpenAI obtained a
lower bound \(u(n)\gg n^{1+\delta}\) for some \(\delta>0\) along a
sequence \cite{OAI1}. Sawin proved that one can take
\(\delta=0.014114\) \cite[Theorem~1]{S}. In September OpenAI proved
\(u(n)=O(n^{4/3-\Delta})\) for some \(\Delta>0\) \cite{OAI2}.
The 17-point diagram in the talk illustrates the counting convention.
Its highlighted Moser spindle consists of two unit rhombi meeting at one
vertex, their short diagonals, and a unit edge joining their far vertices.
Each drawn segment contributes one unordered pair.

\begin{mainthm}\label{thm:main}
Put \(H_0=226(2^{2100}+1)\), \(N_0=\lceil 32\exp(19H_0)\rceil\), and
\begin{equation}\label{eq:delta}
  \Delta=\frac{1}{30600(7200N_0)^4\log(1+N_0^4)}.
\end{equation}
There is an absolute constant \(C>0\) such that
\(u(n)\le Cn^{4/3-\Delta}\) for every integer \(n\ge2\).
\end{mainthm}

All logarithms are natural. To see the size of the saving, the defining
ceiling gives \(32e^{19H_0}\le N_0\le33e^{19H_0}\). Substitution in
\eqref{eq:delta} gives
\[
  2^{2114}\log 2<76H_0<\log(1/\Delta)
  <77H_0+\log 30600<2^{2115}\log 2 .
\]
Thus \(2^{-2^{2115}}<\Delta<2^{-2^{2114}}\). The smaller endpoint
gives the bound stated in the talk since \(n\ge2\). The September
theorem gives no numerical value of its saving with which to compare
\eqref{eq:delta}.

\section{The bounded-variation argument}

We describe the proof through Proposition~4.4 of \cite{OAI2}.
If there were no power saving, some sequence of \(n=t^3\) points would
determine \(t^{4-o(1)}\) unit pairs as \(t\to\infty\).
The incidence reduction gives real algebraic points \(P\), centers \(Q\),
and a unit-incidence graph \(E\subseteq P\times Q\) with
\[
 |P|=t^{1+o(1)},\qquad |Q|=t^{2+o(1)},\qquad |E|=t^{2+o(1)}.
\]
Uniformly, \(d(p)=t^{1+o(1)}\), \(d(q)=t^{o(1)}\), and
\(\min_{q\in Q}d(q)\to\infty\). Earlier forms of the geometric partition
appear in \cite[Corollary~5.11]{katz-silier} and
\cite[\S3]{pach-raz-solymosi}.

Under the uniform edge law \(\omega(p,q)=\mathbf1_E(p,q)/|E|\), the
marginals are \(\mu(p)=d(p)/|E|\) and \(\kappa(q)=d(q)/|E|\).
Choose \(q\sim\kappa\) and two conditionally independent uniform
neighbors \(p,p'\) of \(q\). The wedge law and its endpoint law are
\[
 \Omega(p,q,p')=\frac{\mathbf1_E(p,q)\mathbf1_E(p',q)}{|E|d(q)},
 \qquad \pi(p,p')=\sum_q\Omega(p,q,p').
\]
On the retained graph the mutual information
\(D=\sum_{\pi(p,p')>0}\pi(p,p')\log(\pi(p,p')/(\mu(p)\mu(p')))\)
is \(o(\log t)\) \cite[\S2]{OAI2}.
The geometry enters even here. Two distinct endpoints have at most two
common centers, so their joint mass is at most \(2/(|E|\min_q d(q))\).
The diagonal has total mass \(|Q|/|E|=o(1)\). These bounds keep the
joint entropy within \(o(\log t)\) of the sum of the marginal entropies.

Write \(z=x+iy,\ w=x-iy\); a unit edge \(e=(p,q)\) has jumps
\(d_e=z_p-z_q\) and \(d_e^{-1}=w_p-w_q\).
At each number-field place \(v\) and threshold \(s\), split the edges
according as \(-\log|d_e|_v>s\) or not. Let \(S\) be the sum over places,
with the usual local-degree weights, of the integral over \(s\) of the
smaller of the two edge proportions. Assuming \(S\le S_0\), one may
translate and rescale the coordinates to \(Z,W\), preserving the
unit-edge product and bounding their mean logarithmic heights.

A test chooses a uniform embedding of a common normal number field into
\(\mathbb C\) and records the mesh cell of \((\sigma Z,\sigma W)\), with an overflow
code. A center accepts the cells attained by its neighbors. Since the
alphabet size \(M\) is fixed and \(D=o(\log t)\), the prediction lemma
of \cite[\S3]{OAI2} gives a history \(H\) of \(k=o(\log t)\) tests
with at most \(M^k=t^{o(1)}\) values. First draw \(p\sim\mu\). Given
\(p\), draw \(\widehat p\sim\mu(\cdot\mid H=H(p))\) and three uniform
neighbors \(q_1,q_2,q_3\) of \(p\), independently. At a single fresh
embedding each \(q_j\) accepts \(\widehat p\) with probability \(1-o(1)\).
A union bound gives simultaneous acceptance.

For this short history, choose \(B\ge2\) large and \(0<r<1\) small,
depending only on \(S_0\). With probability at least \(0.99\) for large \(t\),
the sampled coordinates, edge jumps and inverse edge jumps have modulus
at most \(B\),
\(\lvert Z_{\widehat p}-Z_p\rvert\ge r\), and the jumps
\(D_j=Z_p-Z_{q_j}\) differ pairwise by at least \(r\).
Here coordinates denote their embedded values. The height bounds make
upper and lower tail probabilities for nonzero differences arbitrarily
small, uniformly in the fixed history. For the short history the point collision
probability is at most \(|H(P)|\max_p\mu(p)=o(1)\), and each center
collision has probability at most \(1/\min_p d(p)=o(1)\).

Choose \(0<\tau<\min(1,r^4/(128B^5))\) for the mesh width. On a
simultaneous accepted and separated outcome, comparing the candidate
with a genuine neighbor in its cell gives
\(\lvert(Z_{\widehat p}-Z_{q_j})(W_{\widehat p}-W_{q_j})-1\rvert
\le 8B\tau\).
Subtract the exact equation for \(p\) and then take differences between
centers. The resulting two linear equations in
\(Z_{\widehat p}-Z_p,W_{\widehat p}-W_p\) have determinant of modulus
\(\prod_{i<j}|D_i-D_j|/|D_1D_2D_3|\ge r^3/B^3\).
Cramer's rule gives
\(\lvert Z_{\widehat p}-Z_p\rvert\le64B^5\tau/r^3<r/2\), a contradiction.
Thus \(S\to\infty\) on a near-extremal sequence. The slide's three thin
annuli illustrate the argument at the real-plane embedding. At a general
embedding the two coordinates need not be conjugate. The product identity holds
at every embedding.

This identity also explains the role of the norm. For the strictly convex
unit ball \(x^2+|y|\le1\), the vector \((1/2,3/4)\) is a unit vector,
but its Euclidean product \(zw\) equals \(13/16\). Indeed, on the grid
\((i/m,j/m^2)\), \(0\le i\le m,\ 0\le j\le m^2\), each displacement
with index difference \((h,m^2-h^2)\) for \(m/3\le h\le2m/3\) occurs
\((m+1-h)(h^2+1)\gg m^3\) times and is a unit displacement for this
norm. This gives \(\gg m^4\) unordered unit pairs among \(\asymp m^3\)
points, as in Valtr's example \cite{valtr}.

\section{How the argument gives an explicit saving}
\subsection{Counting histories fixes the exponent}
Assume that Theorem~\ref{thm:main} fails. For arbitrarily large
\(t=n^{1/3}\) there are then at least \(t^{4-3\Delta}\) ordered unit
pairs. The quantitative incidence reduction gives a retained graph with
\[
 D+1\le30600\Delta\log t,\qquad
 t\max_p\mu(p)\le 1/100,\qquad \min_p d(p)\ge300 .
\]
Set \(S_0=2^{2100}\). If \(S\le S_0\), normalization gives
\[
 \mathbb E_\mu[h(Z)+h(W)]+\mathbb E_\kappa[h(Z)+h(W)]
 \le226(S+1)\le H_0 ,
\]
where \(h\) is the absolute logarithmic Weil height. Sawin's example
over a CM field illustrates how a short history of embeddings permits
resampling \cite[post 9400]{sawin}. The following test works for the
general number field.

Take \(R=e^{6H_0}\), \(\eta=e^{-7H_0}\), and mesh width at most
\(\eta/(16R)\) on each real coordinate in \([-R,R]\). This takes
\(N_0=\lceil32e^{19H_0}\rceil\) intervals per coordinate, hence
\(M=1+N_0^4\) codes including overflow. A bounded center accepts
overflow and the cells occupied by its genuine neighbors. An overflowing
center accepts every code. On the part \(|Z-Z_q|\ge1\), the identity
\(W-W_q=(Z-Z_q)^{-1}\) shows that differences in \(W\) are bounded by
the corresponding differences in \(Z\). On the remaining part one can
interchange \(Z\) and \(W\). Thus a bounded center accepts at most
\(32N_0^2+1\) codes. The quantitative prediction estimate gives fresh
rejection probability at most \(1/200\) with a history of length
\(k<A(D+1)\), where \(A=(7200N_0)^4\). The choice of \(\Delta\) now
has a precise purpose:
\begin{equation}\label{eq:balance}
 k\log M<A(D+1)\log M
 \le 30600\Delta A\log M\,\log t=\log t .
\end{equation}
There are fewer than \(t\) histories. The probability of resampling
\(p\) itself is therefore at most \(t\max_p\mu(p)\le1/100\).

It remains to turn acceptance into an exact equation. Put
\(F=(Z_{\widehat p}-Z_q)(W_{\widehat p}-W_q)-1\).
In an accepted nonoverflow cell with bounded center, comparison with a
genuine neighbor gives \(|\sigma F|\le\eta/2\). The height inequalities
give \(\mathbb E h(F)\le H_0+3\log2\). For \(F\ne0\),
\(h(F^{-1})=h(F)\), so a fresh embedding seldom makes a nonzero
residual this small. Including overflow, the probability of acceptance
with \(F\ne0\) is at most
\[
 \frac16+\frac{H_0+3\log2}{7H_0}<0.311 .
\]
The first term bounds overflow by Markov's inequality. A repeated
center costs at most \(3/300\). The total failure probability is at most
\(1/100+3/300+3(0.311+1/200)<1\). Thus two distinct points would
have three distinct exact centers. Indeed, with
\(u=Z_{\widehat p}-Z_p\ne0\), \(v=W_{\widehat p}-W_p\), each of the
three distinct jumps \(d=Z_p-Z_q\) would solve \(vd^2+uvd+u=0\).
This is impossible. We conclude that \(S>S_0\).

\subsection{Why the large-variation case preserves the saving}
For two points \(p,p'\) with a common center \(q\), record their direction
ratio \(\lambda=(z_p-z_q)/(z_{p'}-z_q)\). It satisfies
\(\lVert p-p'\rVert^2=2-\lambda-\lambda^{-1}\). Sawin proposed fixing
the alternating products of such ratios around rectangles and studying
which directions can still vary \cite[post 9389]{sawin}. The conditional
grid refinement in \cite{arun} reduces the generic configuration from
\(20\times5\) to \(8\times4\). We use this smaller grid and a height
bound for one of its coordinates. We outline this part of the research
formalization.

Moment estimates on the original common-center walk yield a scale
\(J\ge S/500\) and a dense graph of point pairs with assigned ratios.
Every retained ratio has height at least \(J/3\). The alternating
quotient \(\lambda_{ij}\lambda_{i'j'}/(\lambda_{ij'}\lambda_{i'j})\)
on each complete rectangle has height at most
\(J/(324\cdot 2^{111})\). Freeze these quotients and allow the positions
and ratios to vary subject to their equations. If no tangent vector
moves every ratio to first order, one ratio is stationary on the whole
tangent space: a vector space over \(\mathbb C\) cannot be covered by
finitely many proper linear subspaces. The
coefficient estimate for that ratio uses at most 24 normalized quotients
and gives
\[
 h(\lambda_{ij})\le 2^{2086}
  +2^{113}\frac{24J}{324\cdot2^{111}}
 \le\left(\frac{125}{4096}+\frac{8}{27}\right)J<\frac{J}{3}.
\]
Here \(2^{2086}\le 125J/4096\) follows from \(J\ge2^{2100}/500\).
The lower bound on each individual height excludes the stationary case.

Four fixed anchors then constrain their common neighbors to a curve of
degree at most \(80\). After restricting both sides to curves and
excluding the case of a common unit circle rich in retained points, the
position tuples of compatible \(8\times4\) grids lie in a proper
algebraic subset of the product of the twelve curves. Its defining
polynomial has uniformly bounded degree. Write \(m\) for the number of
points supporting the extracted pair graph. Its restriction to the two
curves retains \(e\gg m^{2-25/2000}\) edges on sides of sizes \(a,b\le m\).
Count ordered grids with repetitions. If \(d(\mathbf p)\) is the common
degree of an ordered eight-tuple on the left, the power-mean inequality
applied twice gives
\[
 \sum_{\mathbf p}d(\mathbf p)^4
 \ge a^{-24}\Bigl(\sum_{r\ {\rm on\ the\ right}}d(r)^8\Bigr)^4
 \ge \frac{e^{32}}{a^{24}b^{28}}
 \gg m^{11+3/5}.
\]
Polynomial zero counting on the product of twelve curves gives
\(O(m^{11})\) for the same tuples. These bounds contradict one another for sufficiently large
\(m\). The polynomial degree affects how large \(m\) must be. It costs
no further power of \(n\) in \eqref{eq:balance}.

\section{Work during the 40 hours and formal verification}
We worked through the September proof up to Proposition~4.4 and
identified the Euclidean identity used in the three-center argument.
We first stated the exact exponent and its high-level dependencies in
Lean. Checking those statements led us to keep the original walk in the
moment estimates and to bound the height of one stationary ratio.
We then checked the formal conclusion against the literal count of
unordered pairs in finite subsets of \(\mathbb R^2\).

The following are the declarations used here, all in the namespace
\texttt{PlanarUnitDistances}. Their hypotheses specify where each
intermediate estimate applies.
\begin{itemize}
\item \texttt{ResearchPhysical.original\_walk\_moments} gives the three
moment bounds and the two endpoint estimates for a sufficiently large
retained incidence graph.
\item \texttt{GridLocus.proper\_grid\_locus} gives the separating-polynomial
statement for the admissible ordered pair of curves.
\item \texttt{IntegralGrid.stationary\_height\_unconditional} gives the
stationary-coordinate height estimate with cap \(2^{2086}+2^{113}s\),
where \(s\) is the sum of the relevant quotient heights.
\item \texttt{exactGap\_log\_bounds} gives the two middle bounds on
\(\log(1/\Delta)\) in Section~1. The outer bounds there are numerical.
\item \texttt{IntegralGrid.exact\_unconditional} proves
Theorem~\ref{thm:main} with its displayed \(\Delta\), no hypotheses,
and \(n\ge2\). One \(C>0\) works for every \(n\) in that range.
\end{itemize}

The research project builds with Lean and Mathlib v4.35.0. Parts of its
incidence, counting, specialization and probability arguments adapt
OpenAI's published Lean sources \cite{OAI2}. We claim the explicit
quantitative refinement and report the kernel check with this attribution.
The project does not qualify as an independent competition formalization
of the official OpenAI theorem. The official declaration asks for
an unspecified exponent and includes \(n=0,1\).

\begin{center}
\begin{tabular}{@{}lr@{}}
\toprule
Quantity in the dependency chain of Theorem~A & Count\\
\midrule
Proof modules & 312\\
Lines containing Lean code in those modules & 31,592\\
Source declarations, including adapted code & 2,536\\
Project declarations replayed, including generated declarations & 3,728\\
Direct dependency pairs between project declarations & 19,099\\
Transitive external declarations inspected & 77,801\\
Admitted proofs in the dependency chain & 0\\
Additional project axioms & 0\\
\bottomrule
\end{tabular}
\end{center}

The line count excludes blank lines and lines consisting only of
comments. A fresh kernel replay started at the named theorem and checked
all project declarations on which its type or proof depends. The only
axioms reported were \texttt{propext}, \texttt{Classical.choice}, and
\texttt{Quot.sound}.

\subsection*{AI use and further work}
We used OpenAI Codex to develop Lean proofs, compare constants and
quantifiers across the manuscript and code, and draft the exposition.
We checked the sampling formulas against the September paper, the
displayed exponent against its Lean definition, and the final declaration
by kernel replay. Our detailed human reading during the event reached
Proposition~4.4. The next step is to work through the large-variation
rigidity argument and its formal proof in full, and to write a
self-contained human account of the explicit saving.

\begin{thebibliography}{10}
\bibitem{E} P. Erd\H{o}s, \href{https://doi.org/10.1080/00029890.1946.11991674}
{On sets of distances of \(n\) points}, \emph{Amer. Math. Monthly}
\textbf{53} (1946), 248--250.

\bibitem{katz-silier} N. H. Katz and O. Silier,
\href{https://arxiv.org/abs/2303.17186v2}{Structure of cell decompositions
in extremal Szemer\'edi--Trotter examples}, arXiv:2303.17186,
2023, rev. 2026.

\bibitem{OAI1} OpenAI,
\href{https://cdn.openai.com/pdf/74c24085-19b0-4534-9c90-465b8e29ad73/unit-distance-proof.pdf}
{Planar point sets with many unit distances}, preprint, May 2026.

\bibitem{OAI2} OpenAI,
\href{https://github.com/openai/math/blob/adc7f1241b42e322a6451854ab7e4b4c146bf78a/preprints/A-power-saving-for-planar-unit-distances-September-23-2026/paper.pdf}
{A power saving for planar unit distances}, preprint, 23 September 2026,
with \href{https://github.com/openai/math/tree/adc7f1241b42e322a6451854ab7e4b4c146bf78a/lean/OAI/Geometry/UnitDistances}
{accompanying Lean sources}.

\bibitem{pach-raz-solymosi} J. Pach, O. E. Raz and J. Solymosi,
\href{https://doi.org/10.4230/LIPIcs.SoCG.2026.83}{Erd\H{o}s's unit
distance problem and rigidity}, in \emph{42nd International Symposium
on Computational Geometry}, LIPIcs \textbf{367} (2026), article 83.

\bibitem{arun} \texttt{rohanarun} (repository maintainer),
\href{https://github.com/rohanarun/unit-distance-grid-refinement/tree/4e7b3af0a46a7127a57019a0ab6d76fa2b55fe2e}
{A smaller generic grid in the unit-distance argument}, research draft
prepared with Codex and Claude, 10 October 2026, revised 13 October 2026.

\bibitem{S} W. Sawin, \href{https://arxiv.org/abs/2605.20579}
{An explicit lower bound for the unit distance problem},
arXiv:2605.20579, 2026.

\bibitem{sawin} W. Sawin, comments on T. Bloom,
\emph{Erd\H{o}s problems in the OpenAI proofs}, Erd\H{o}s Problems,
\href{https://www.erdosproblems.com/forum/thread/blog:10\#post-9389}
{post 9389, 25 October 2026}, and
\href{https://www.erdosproblems.com/forum/thread/blog:10\#post-9400}
{post 9400, 1 November 2026}.

\bibitem{SST} J. Spencer, E. Szemer\'edi and W. T. Trotter,
Unit distances in the Euclidean plane, in \emph{Graph Theory and
Combinatorics} (B. Bollob\'as, ed.), Academic Press, 1984, 293--303.

\bibitem{valtr} P. Valtr,
\href{https://kam.mff.cuni.cz/~valtr/n.pdf}{Strictly convex norms allowing
many unit distances and related touching questions}, manuscript, \S2.1.
\end{thebibliography}
\end{document}
